
\section{Implementation and Experimental Design}
\label{sec:06}
\label{sec:implementation-experiments}

\subsection{Implementation Configuration}
\label{subsec:implementation-profile}



VoteChain-ZK extends the established biometric voting architecture
with a cryptographic ballot-processing protocol
\citep{rifat2026biometric}. The implementation separates voter
authentication from encrypted ballot construction. A voter-facing
web client manages candidate selection and initiates biometric
authorization. The credential authorization service authenticates
the voter and releases the protected voting credential. A
client-side privacy worker then prepares the encrypted ballot
and generates its validity proof.
The web client is implemented using React and TypeScript.
The authentication and election management services use Node.js.
The privacy worker performs credential recovery, ballot encryption,
witness generation, and local proof verification. Face and fingerprint
recognition are provided by the existing authentication components.
The cryptographic experiments do not repeat the biometric
recognition evaluation reported in the earlier study.

The voting relation is implemented in Circom 2. The selected
circuit uses a fixed Merkle depth of 17 and represents each
ballot through 16 encrypted candidate slots. Witness generation
and Groth16 proving are performed using \texttt{snarkjs}.
A generated Solidity verifier checks the submitted proof.
The \texttt{VoteChainDirect} contract enforces election-state
conditions, rejects previously used nullifiers, and updates
the encrypted tally after successful verification.

The system also implements threshold-based result recovery.
Five trustees participate under a reconstruction threshold
of three. The threshold protocol verifies distributed
key-generation information, partial decryptions, and
their accompanying consistency proofs. Valid contributions
are combined to recover aggregate candidate counts.
These cryptographic checks are performed outside the
voting contract. The contract subsequently records the
result values and finalization-transcript hash after
checking its specified state and arithmetic conditions.
The evaluated cryptographic artifacts use a development-stage
Groth16 trusted setup. Blockchain execution takes place
within a controlled local environment. The experimental
design therefore measures the implemented protocol under
research conditions rather than a deployed public election.


% Implementation evidence at pinned revision:
% 17c7453c50793643b590fb5beac6ad77b9aa53ce
% crypto/config/wp5-application-profile.json
% circuits/VoteChainVoteDirect.circom
% votechain-frontend-ts/src/privacy/wp5Privacy.worker.ts
% contracts/wp4-foundry/src/VoteChainWP4.sol
% crypto/src/wp6-threshold.mjs
% votechain-central-server/services/elections/electionFinalizationService.js

\subsection{Experimental Environment and Measurement Scope}
\label{subsec:experimental-environment}



The experiments were conducted on a Windows-based AMD64
computer with 16 logical processors and approximately
10.6~GB of physical memory. The recorded software
environment included Python 3.10.11, Node.js, Circom,
\texttt{snarkjs}, Foundry, and Anvil.
Blockchain measurements were performed using a local
Anvil network. Transaction gas consumption was obtained
from successful encrypted-ballot submissions. Transaction
latency was measured from the beginning of submission
to the return of the corresponding receipt. This interval
captures local transaction execution and receipt retrieval.
It does not include biometric authentication, credential
authorization, or proof preparation.

Cryptographic operations were measured independently
from blockchain execution. The measured stages include
eligible-identity materialization, Merkle-tree preparation,
circuit witness generation, Groth16 proof generation,
and local proof verification.
The evaluation comprises four experimental campaigns.
Each campaign examines a different computational
condition. Accepted-ballot transactions are evaluated
separately from fresh-chain execution. Candidate-count
and eligible-set experiments assess the effects of
their respective configuration parameters.

\subsection{Performance Campaign Design}
\label{subsec:experimental-campaigns}



Table~\ref{tab:implementation-campaigns} summarizes
the completed experimental workloads. The campaigns
use repeated measurements and exclude warm-up
observations from their statistical summaries.

\begin{table}[t]
\centering
\caption{Experimental configurations and measured sample sizes.}
\label{tab:implementation-campaigns}
\small
\begin{tabular}{@{}p{0.23\linewidth}p{0.44\linewidth}p{0.25\linewidth}@{}}
\toprule
\textbf{Campaign} &
\textbf{Experimental configuration} &
\textbf{Measured observations} \\
\midrule
Accepted ballots
& 25 independent elections, with five elections at each
size of 25, 50, 75, 100, and 150 accepted ballots
& 2,000 transactions \\
\addlinespace
Fresh-chain first ballot
& First accepted private ballot following a fresh
local-chain setup
& 20 trials \\
\addlinespace
Candidate count
& Active candidate counts of 2, 4, 8, and 16
with a fixed 16-slot representation
& 20 trials per count \\
\addlinespace
Eligible set
& Eligible populations of 10, 100, 1,000, and
10,000 with Merkle depth fixed at 17
& 20 trials per population \\
\bottomrule
\end{tabular}
\end{table}

\paragraph{Accepted-ballot transactions.}

The accepted-ballot experiment examines transaction
costs across successive votes within an election.
Five independent elections were evaluated at each
of five election sizes. These sizes were 25, 50,
75, 100, and 150 accepted ballots. The resulting
25 elections produced 2,000 measured transactions.
An additional 400 warm-up transactions were excluded
from the statistical analysis. Gas consumption and
local submission-to-receipt latency were measured
for each accepted transaction. The experiment
captures variation across the recorded election
workloads rather than considering only the first
vote after contract deployment.

\paragraph{Fresh-chain execution.}

The fresh-chain experiment isolates the first
accepted private ballot following local-chain
initialization. Each iteration begins with a
fresh execution environment. The measurement
interval starts at transaction submission and
ends when the corresponding receipt is returned.
Contract preparation and post-receipt checks
remain outside this interval. The experiment
contains 20 measured trials. Five additional
warm-up trials were excluded from the final
statistics.
This experiment provides a separate measurement
of first-ballot execution under a fresh contract
state. Its observations are not pooled with
the successive-ballot campaign.

\paragraph{Candidate-count sensitivity.}

The candidate-count experiment evaluates four
configurations containing 2, 4, 8, and 16
active candidates. Each configuration contributes
20 measured transactions. Five warm-up trials
per configuration were excluded. Candidate
positions were balanced across the recorded
trial sequence.
The cryptographic circuit maintains the same
16-slot representation in all configurations.
Consequently, the circuit dimensions and
ciphertext structure remain fixed while
the active candidate count changes.
This experiment evaluates transaction behavior
under different active candidate counts. It
does not measure the cost of increasing the
number of circuit slots or deploying circuits
of different widths.

\paragraph{Eligible-set sensitivity.}

The eligible-set experiment evaluates synthetic
populations of 10, 100, 1,000, and 10,000
registered commitments. The selected voter
occupies the first indexed position in each
eligible-set fixture.
The Merkle depth remains fixed at 17,
corresponding to a nominal capacity of
131,072 leaves. The ballot circuit also
retains its fixed 16-slot representation.
Each population contributes 20 measured runs
following five warm-up runs. The execution
order varies across population sizes to
distribute their positions within the
measurement sequence.

Five computational stages are measured separately.
The first stage materializes eligible identity
leaves. The second constructs the Merkle snapshot,
root, and target authentication path.
The remaining stages measure circuit witness
generation, Groth16 proof generation, and
local Groth16 proof verification.
The campaign does not submit blockchain
transactions. Its verification measurements
therefore represent off-chain proof verification
rather than Solidity verifier execution.
Because the Merkle depth remains fixed,
this experiment measures the effect of
eligible-set population on the specified
operations. It does not measure the
performance of circuits with increasing
Merkle depth.


% Experimental evidence:
% experiments/wp8/20261002_232814/wp8e-c7-final-statistics/
%   C7_FINAL_STATISTICAL_RECORD.txt
% experiments/wp8/20261003_020418/
%   wp8e-c8b-r4-final-microbenchmark-statistics/
%   C8B_R4_FINAL_STATISTICS.json
% experiments/wp8/20261003_212713/
%   wp8e-c8c-r7f-final-measured-statistics/
%   C8C_R7F_STATISTICAL_VALIDATION.json
% experiments/wp8/20261004_014001/
%   wp8e-c8d-r8-final-100-observation-campaign/
%   C8D_R8_FINAL_CAMPAIGN_VALIDATION.json

\subsection{Statistical Analysis and Validity Checks}
\label{subsec:experimental-analysis}



Statistical analysis was performed using the measured
observations from each campaign. Warm-up trials were
excluded according to the experimental protocol.
Pilot and aborted campaign observations were not
combined with the finalized measurements. No measured
outliers were removed from the statistical analyses.
Performance metrics are summarized using sample size,
mean, median, minimum, maximum, sample standard
deviation, and 95th percentile. The sample standard
deviation uses the denominator $n-1$. The 95th
percentile is calculated using the nearest-rank
method.

These statistics describe the observed computational
behavior under the evaluated configurations.
The candidate-count campaign uses descriptive
comparison rather than inferential regression.
The fixed circuit dimensions also define the
scope of the eligible-set measurements.
Cryptographic validity checks accompanied the
eligible-set experiment. Generated witnesses
were checked against their expected constraints.
The corresponding Groth16 proofs were verified
locally. The measured observations also passed
the public-signal consistency checks.

Functional tests additionally examine rejection
conditions and protocol correctness. These
include malformed ballot witnesses, repeated
nullifiers, and threshold-verification conditions.
Such tests are distinct from the positive
performance measurements.
Threshold cryptographic operations were also
examined through functional and integration
checks. However, the four quantitative
campaigns do not provide a repeated timing
estimate for complete threshold finalization.
These measurements are therefore not combined
with accepted-ballot or proof-generation
performance statistics.
The following section presents the measured
results. Their computational and operational
implications are examined in the subsequent
discussion.

